Based on the findings, we present three suggestions for future research, ranging from very related to our work to
a completely novel approach.

In \autoref{ch:ramp_model} we apply a ramp loss based on the distance of a point to a ruleset.
The ramp loss model poses various difficulties.
These include undesirable behavior of the additive variant (\autoref{sec:ramp-examples}),
slow solve times for the pricing problems (even after implementing cutting planes described in \autoref{subsec:ramp-cutting-planes}),
and a complicated pricing problem for the non-additive variant (\autoref{sec:ramp-nonadditive}).
It is arguably difficult to work with and solve.
An alternative strategy is to define an \textit{importance score}
\[
    \chi_i = \left[\frac{1}{\beta}|(\text{PD-L1 expression})_i - \alpha|\right]_0^1,
\]
where $[\cdot]$ is the clipping operator, $\alpha$ is the gene expression threshold and $\beta$ is the margin width.
Then, consider
\[
    \sum_{i \in \P}\chi_i \zeta_i + \sum_{i \in \N} \chi_i \zeta_i
\]
as the loss function in the original master IP \eqref{eq:BM}.
We may also consider a Hamming loss equivalent.
This strategy effectively means that samples whose gene expression is closer to the target threshold are weighed less heavily.
This is a much simpler adjustment of the boolean model or threshold model that is similar conceptually.
In fact, it may avoid many of the artifacts observed in the ramp loss model and difficulties encountered in its implementation.

Another suggestion is to re-evaluate use of DecoupleR as unquestioned input data.
DecoupleR itself is an estimation method.
In particular, it is one that assumes that TF affect target genes individually.
Rather than treating TF activity estimation as data pre-processing,
a suggestion is to unite TF activity score estimation and the search for TF combinations into a single, novel method.
Perhaps a united approach, one that assumes combinatorial TF regulation from raw data onwards, can yield better results.

Finally, it is suggested to use statistical methods to search for statistically significant concurrent TF activity,
rather than attempt to predict PD-L1 expression using discrete, deterministic classifiers.
An example of such a statistical method, albeit for a different use case, is LAMP~\cite{terada_statistical_2013}.